% ============================================================
% Unconditional Hyperbolicity of the Critical Laguerre Family:
% An Elementary and Formally Verified Proof
%
% arXiv-ready draft v1 — 2026-06-12.
% Title and author block are placeholders pending Michael's decision
% (see ACKNOWLEDGMENTS / AUTHORSHIP note near \author).
% Compile: pdflatex (no external .bib; thebibliography inline).
% ============================================================
\documentclass[11pt]{amsart}

\usepackage[T1]{fontenc}
\usepackage[utf8]{inputenc}
\usepackage{amsmath,amssymb,amsthm}
\usepackage{mathtools}
\usepackage{microtype}
\usepackage[hidelinks]{hyperref}
\raggedbottom

\newtheorem{theorem}{Theorem}[section]
\newtheorem{lemma}[theorem]{Lemma}
\newtheorem{proposition}[theorem]{Proposition}
\newtheorem{corollary}[theorem]{Corollary}
\theoremstyle{definition}
\newtheorem{definition}[theorem]{Definition}
\theoremstyle{remark}
\newtheorem{remark}[theorem]{Remark}

\DeclareMathOperator{\sgn}{sgn}
\newcommand{\dd}{\,\mathrm{d}}
\newcommand{\e}{\mathrm{e}}
\newcommand{\R}{\mathbb{R}}
\newcommand{\C}{\mathbb{C}}
\newcommand{\N}{\mathbb{N}}
\newcommand{\E}{\mathbb{E}}
\newcommand{\Prob}{\mathbb{P}}
\newcommand{\He}{\mathrm{He}}
\newcommand{\hypF}[3]{{}_1F_1\!\left(#1;\,#2;\,#3\right)}

\begin{document}

\title[Hyperbolicity of the critical Laguerre family]{Unconditional
hyperbolicity of the critical Laguerre family:\\ an elementary and
formally verified proof}

% AUTHORSHIP NOTE (to be settled by M.~Haufschild before submission):
% author list, affiliation, and the framing of the AI contribution are
% the repository owner's call. The Lean formalization was produced
% jointly by two AI research agents ("Fable" = Claude, and GPT) under
% the author's direction; see Section 6 and the Acknowledgments.
\author{Michael Haufschild}
\email{haufschildmichael@gmail.com}

\date{June 12, 2026}

\begin{abstract}
For $d\ge 1$ let
$C_d(w)=\hypF{-d}{\tfrac12}{w^2/4d}$ --- up to normalization the
Laguerre polynomial $L_d^{(-1/2)}$ at its critical scaling --- and let
\[
\Psi_d(w)\;=\;\sum_{k=0}^{d}(-1)^k\,
\frac{d^{\underline{k}}\,M_k}{d^{\,k}\,(2k)!}\,w^{2k},
\qquad
M_k=\frac{(2k)!}{4^k\,k!}\,\e^{\gamma k-S_1(k)},
\]
where $d^{\underline{k}}$ is the falling factorial,
$\gamma$ is Euler's constant and $S_1(k)=\sum_{j=1}^{k}H_{j-1}$ is the
partial sum of harmonic numbers. We prove that for every $d\ge 1$ all
$2d$ zeros of $\Psi_d$ are real and simple. The numbers $M_k$ are
shown to be the moment sequence of an explicit compound-Poisson random
variable $U\in(0,1]$ with atom $\Prob(U=1)=\sqrt{2/\e}$, so that
$\Psi_d(w)=\E\,C_d(\sqrt{U}\,w)$: hyperbolicity of the whole family is
an exact averaging, not an asymptotic, phenomenon. The proof is
elementary --- real-variable throughout, using only the Kummer
differential equation, an energy monotonicity along the real axis, and
a sign-alternation count at the critical points of $C_d$ --- and it has
been formally verified: the theorem is proved in Lean~4 against the
mathlib library with zero \texttt{sorry}s, depending only on the three
standard axioms. To our knowledge $\Psi_d$ is, in this generality, a
new explicit hyperbolic deformation of the Laguerre--Hermite family.
\end{abstract}

\maketitle

\section{Introduction}\label{sec:intro}

A real polynomial is \emph{hyperbolic} if all of its complex zeros are
real. The hyperbolicity of the classical orthogonal families ---
Hermite, Laguerre, Jacobi --- is classical
\cite[Thm.~3.3.1]{szego}, and the linear operations preserving
hyperbolicity were characterized by P\'olya and Schur
\cite{polya-schur}. Deciding hyperbolicity of a \emph{specific}
explicitly given family that is not covered by such preserver theorems
remains, in general, a genuinely hard problem: one must either locate
all the zeros or manufacture enough sign changes.

This paper proves hyperbolicity, with simple zeros, for the following
family. Throughout, $d^{\underline{k}}=d(d-1)\cdots(d-k+1)$ denotes the
falling factorial, $H_n=\sum_{j=1}^n 1/j$ the harmonic numbers
($H_0=0$), $\gamma$ Euler's constant, and
\begin{equation}\label{eq:S1}
S_1(k)\;=\;\sum_{j=1}^{k}H_{j-1}\;=\;kH_{k-1}-(k-1)\qquad(k\ge 1),
\qquad S_1(0)=0 .
\end{equation}

\begin{definition}\label{def:family}
For integers $k\ge 0$ and $d\ge 1$ set
\begin{equation}\label{eq:Mk}
M_k\;=\;\frac{(2k)!}{4^k\,k!}\,\e^{\gamma k-S_1(k)},
\end{equation}
\begin{equation}\label{eq:CdPsid}
\begin{aligned}
C_d(w)\;&=\;\sum_{k=0}^{d}(-1)^k\,
\frac{d^{\underline{k}}}{d^{\,k}(2k)!}\,w^{2k}
\;=\;\hypF{-d}{\tfrac12}{\frac{w^2}{4d}},\\
\Psi_d(w)\;&=\;\sum_{k=0}^{d}(-1)^k\,
\frac{d^{\underline{k}}\,M_k}{d^{\,k}(2k)!}\,w^{2k}.
\end{aligned}
\end{equation}
\end{definition}

By Kummer's transformation of the confluent hypergeometric polynomial,
$C_d$ is the Laguerre polynomial $L_d^{(-1/2)}$ at the argument scaling
$w^2/4d$ that places its zeros at unit mean density --- we call
$\{C_d\}_{d \ge 1}$ the \emph{critical Laguerre family}, and
$\Psi_d$ its \emph{moment deformation} by the sequence
\eqref{eq:Mk}. The deformation is of genuinely infinite order: the
sequence $(M_k)$ is strictly decreasing
(the ratio test below reduces this to $\log(1+x)<x$), from $M_0=1$
through $M_1=\e^{\gamma}/2=0.89053\ldots$ toward its limit
$p:=\sqrt{2/\e}=0.85776\ldots$, so every coefficient of $C_d$ is
damped by a different factor. Nothing about \eqref{eq:Mk} makes
$(M_k)$ visibly a multiplier sequence in the sense of
P\'olya--Schur, and we do not know a preserver-theoretic proof of
Theorem~\ref{thm:main}.

\begin{theorem}[Main theorem]\label{thm:main}
For every integer $d\ge 1$, all $2d$ zeros of $\Psi_d$ are real and
simple.
\end{theorem}

The proof has two ingredients, both of independent interest.

\emph{First ingredient: the moments are a compound-Poisson law.} The
sequence \eqref{eq:Mk} is exactly the moment sequence of an explicit
random variable: let $N$ be a Poisson point process on $(0,\infty)$
with the finite intensity measure
\[
\lambda(\dd t)\;=\;\eta(t)\,\frac{\e^{-t}}{1-\e^{-t}}\,\dd t,
\qquad
\eta(t)\;=\;\frac{\e^{t/2}}{t}-\frac{1}{1-\e^{-t}}\;>\;0 ,
\]
and set $U=\e^{-\sum_j T_j}$, the product over the points of $N$. Then
(Theorem~\ref{thm:moments})
\[
\E\,[U^k]\;=\;M_k \quad (k\ge 0),
\qquad
\Prob(U=1)\;=\;p\;=\;\sqrt{2/\e}\,.
\]
Consequently, writing $\nu$ for the law of $v=\sqrt{U}$ and
$\mu=\nu-p\,\delta_1$ for its non-atomic part,
\begin{equation}\label{eq:decomp}
\Psi_d(w)\;=\;\E\,C_d\bigl(\sqrt{U}\,w\bigr)
\;=\;p\,C_d(w)\;+\;\int_{(0,1)}C_d(vw)\,\dd\mu(v),
\end{equation}
with $\mu\ge 0$ and $\mu\bigl((0,1)\bigr)=1-p$.
The identity \eqref{eq:decomp} is exact, coefficient by coefficient,
at every finite $d$; no limit is involved. Nor is the intensity
$\lambda$ an ansatz: Section~\ref{sec:moments} opens by deriving it
from the moment ratios of \eqref{eq:Mk}, through the classical
integral representations of the digamma function.

\emph{Second ingredient: a one-line energy bound at the critical
points.} Writing $E=C_d^2+C_d'^2$, the Kummer differential equation
$C_d''-(w/2d)\,C_d'+C_d=0$ gives $E'(x)=(x/d)\,C_d'(x)^2\ge 0$ on
$x\ge 0$; since $E(0)=1$, every critical point $c$ of $C_d$ satisfies
$|C_d(c)|\ge 1$ and, crucially, the \emph{whole segment} $[0,c]$ is
dominated: $|C_d(vc)|\le |C_d(c)|$ for $0\le v\le 1$. Feeding this
into \eqref{eq:decomp} and using only $p>\tfrac12$ shows that
$\Psi_d$ \emph{copies the sign pattern} of $C_d$ at the critical
points of $C_d$. The classical interlacing data of the Laguerre zeros
then forces $d$ sign changes of $\Psi_d$ on $(0,\infty)$, which (with
evenness and the degree count) pins all $2d$ zeros to the real axis,
each simple.

No complex analysis enters: no Rouch\'e contour, no
Hermite--Biehler theory, no saddle points. The argument runs entirely
on the real line, which is also what made a complete formal
verification practical.

\emph{Formal verification.} Theorem~\ref{thm:main} has been proved in
Lean~4 \cite{lean4} against the mathlib library \cite{mathlib}: the
statement
\begin{center}
\texttt{theorem theorem\_M (d :\ $\N$) (hd :\ 1 $\le$ d) :\\
$\forall$ z $\in$ ((Psi d).map (algebraMap $\R$ $\C$)).roots, z.im = 0}
\end{center}
compiles with zero \texttt{sorry}s and depends only on the three
standard axioms of the Lean/mathlib universe
(\texttt{propext}, \texttt{Classical.choice}, \texttt{Quot.sound}).
Section~\ref{sec:lean} describes the formalization; the source is
distributed with this paper.

\subsection*{Motivation}
The identity \eqref{eq:decomp} exhibits $\Psi_d$ as an exact average
of rescaled copies of $C_d$, and Theorem~\ref{thm:main} says that this
particular averaging preserves hyperbolicity at every degree. Read as
a statement about the sequence $(M_k)$, it is a concrete instance of a
natural question in the tradition of P\'olya--Schur preserver theory:
\emph{which moment sequences of random variables in $[0,1]$ act
coefficientwise on the critical Laguerre family without destroying
hyperbolicity?} The compound-Poisson mechanism of
Section~\ref{sec:moments} --- an atom of mass $p>\tfrac12$ at the
identity scale plus a positive residual --- combined with the energy
bound of Section~\ref{sec:laguerre} yields a sufficient criterion of
striking simplicity, and the sequence \eqref{eq:Mk}, with its
digamma--Binet moment structure and atom exactly $\sqrt{2/\e}$, is a
distinguished instance in which the criterion is endpoint-sharp
(Remark~\ref{rem:sharp}). Beyond that, the theorem stands on its own:
an explicit, nontrivially deformed Laguerre family that remains
hyperbolic at every degree, with an elementary, machine-checked proof
and an exact probabilistic mechanism behind it.

\subsection*{Organization} Section~\ref{sec:moments} constructs the
compound-Poisson law and proves $\E[U^k]=M_k$.
Section~\ref{sec:laguerre} collects the classical facts about $C_d$
and proves the energy lemma. Section~\ref{sec:proof} proves
Theorem~\ref{thm:main}. Section~\ref{sec:lean} documents the formal
verification. Section~\ref{sec:remarks} closes with remarks on
sharpness and an open question.

\section{The compound-Poisson law behind the moments}\label{sec:moments}

\subsection*{Reverse-engineering the intensity}
Before verifying anything, we explain where $\lambda$ comes from:
it is forced by the moments, and the harmonic numbers in
\eqref{eq:Mk} dictate its shape through the digamma function. The
ratios of consecutive moments are elementary --- by the telescoping
$S_1(k)-S_1(k-1)=H_{k-1}$ of \eqref{eq:S1},
\[
\frac{M_k}{M_{k-1}}
=\frac{2k(2k-1)}{4k}\,\e^{\gamma-H_{k-1}}
=\Bigl(k-\tfrac12\Bigr)\,\e^{\gamma-H_{k-1}} .
\]
Now look for a random variable $U=\e^{-Y}\in(0,1]$ with these moments,
taking $Y\ge0$ infinitely divisible --- the simplest class with fully
explicit Laplace transforms: $\log\E[\e^{-kY}]
=\int_0^\infty(\e^{-kt}-1)\,\lambda(\dd t)$ for an intensity measure
$\lambda$ to be found. Consecutive log-moment differences are then
$\int(\e^{-kt}-\e^{-(k-1)t})\,\lambda(\dd t)$, and the ansatz
$\lambda(\dd t)=\eta(t)\,\frac{\e^{-t}}{1-\e^{-t}}\,\dd t$ makes the
difference collapse to a single Laplace transform, since
$(\e^{-kt}-\e^{-(k-1)t})\,\frac{\e^{-t}}{1-\e^{-t}}=-\e^{-kt}$.
Matching the moment ratios above therefore amounts to solving
\[
\int_0^\infty \e^{-kt}\,\eta(t)\,\dd t
\;=\;\psi(k)-\log\Bigl(k-\tfrac12\Bigr),
\qquad k=1,2,\dots,
\]
where the harmonic numbers have become the digamma function through
$\psi(k)=H_{k-1}-\gamma$. This transform pair is classical: Gauss's
integral represents $\psi$, Frullani's integral represents the
logarithm (the half-integer shift $k-\tfrac12$ producing the factor
$\e^{t/2}$), and subtracting the two integrands term by term reads off
\[
\eta(t)\;=\;\frac{\e^{t/2}}{t}-\frac{1}{1-\e^{-t}} .
\]
Everything to this point is bookkeeping. The one substantive fact ---
the fact that makes the random variable exist and the whole method
work --- is that this particular difference of classical integrands is
\emph{positive} and integrable against
$\frac{\e^{-t}}{1-\e^{-t}}\,\dd t$. That is Lemma~\ref{lem:eta}, and
the rest of the section re-runs the derivation forwards, rigorously.

\begin{lemma}\label{lem:eta}
For $t>0$ let
\(
\eta(t)=\dfrac{\e^{t/2}}{t}-\dfrac{1}{1-\e^{-t}}.
\)
Then
\begin{gather*}
\eta(t)\;=\;\frac{2\sinh(t/2)-t}{t\,(1-\e^{-t})}\;>\;0,
\qquad
\eta(t)=\frac{t}{24}+O(t^2)\ \ (t\downarrow 0),\\
\eta(t)\,\frac{\e^{-t}}{1-\e^{-t}}=O\!\bigl(t^{-1}\e^{-t/2}\bigr)\ \ (t\to\infty).
\end{gather*}
In particular
\(
\lambda(\dd t):=\eta(t)\,\frac{\e^{-t}}{1-\e^{-t}}\,\dd t
\)
is a finite positive measure on $(0,\infty)$.
\end{lemma}

\begin{proof}
Multiplying $\eta$ by $t(1-\e^{-t})>0$ gives
$\e^{t/2}(1-\e^{-t})-t=\e^{t/2}-\e^{-t/2}-t=2\sinh(t/2)-t>0$ for
$t>0$, since $\sinh u>u$ for $u>0$. The expansion at $0$ is immediate
from $2\sinh(t/2)-t=t^3/24+O(t^5)$ and $t(1-\e^{-t})=t^2+O(t^3)$; the
decay at infinity from $2\sinh(t/2)\le \e^{t/2}$. Both endpoints are
integrable against $\e^{-t}/(1-\e^{-t})\,\dd t \sim \dd t/t$ at $0$ and
$\sim\e^{-t}\dd t$ at $\infty$, whence $\lambda$ is finite.
\end{proof}

\begin{definition}\label{def:U}
Let $N=\{T_j\}$ be a Poisson point process on $(0,\infty)$ with
intensity $\lambda$ (a.s.\ finitely many points, since
$\lambda$ is finite), and set
\[
Y\;=\;\sum_j T_j\;\in\;[0,\infty),
\qquad
U\;=\;\e^{-Y}\;\in\;(0,1] .
\]
\end{definition}

The variable $U$ is a compound-Poisson exponential functional; its
atom at $1$ is the no-point event,
$\Prob(U=1)=\e^{-\lambda((0,\infty))}>0$. By Campbell's theorem
(the Laplace functional of a Poisson process,
\cite[Sec.~3.2]{kingman}),
\begin{equation}\label{eq:CP}
\E\,[U^k]\;=\;\E\,\e^{-kY}
\;=\;\exp\!\Bigl(\int_0^\infty\bigl(\e^{-kt}-1\bigr)\lambda(\dd t)\Bigr),
\qquad k\ge 0 .
\end{equation}

The transform of $\eta$ is a Binet-type integral. We only need it at
the positive integers, where it is elementary:

\begin{lemma}[integer Binet identity]\label{lem:binet}
For every integer $k\ge 1$,
\[
\int_0^\infty \e^{-kt}\,\eta(t)\,\dd t
\;=\;H_{k-1}-\gamma-\log\Bigl(k-\tfrac12\Bigr).
\]
\end{lemma}

\begin{proof}
Split $\eta(t)\e^{-kt}=\bigl(\e^{-(k-1/2)t}-\e^{-t}\bigr)/t
+\bigl(\e^{-t}/t-\e^{-kt}/(1-\e^{-t})\bigr)$. Each part is absolutely
integrable: the first is bounded near $0$ (numerator $=O(t)$) and
decays exponentially; the second is Gauss's integrand for the digamma
function. Frullani's integral \cite[Sec.~12.16]{ww} gives
\[
\int_0^\infty\frac{\e^{-(k-1/2)t}-\e^{-t}}{t}\,\dd t
=-\log\Bigl(k-\tfrac12\Bigr),
\]
and Gauss's formula \cite[Sec.~12.3]{ww}
\[
\int_0^\infty\Bigl(\frac{\e^{-t}}{t}-\frac{\e^{-kt}}{1-\e^{-t}}\Bigr)\dd t
=\psi(k)=H_{k-1}-\gamma
\]
for the digamma function $\psi=\Gamma'/\Gamma$ at the positive integer
$k$. Adding the two evaluations gives the claim.
\end{proof}

\begin{theorem}[moment identification]\label{thm:moments}
$\E\,[U^k]=M_k$ for every integer $k\ge 0$, and
$\Prob(U=1)=p=\sqrt{2/\e}$. Consequently, with $\nu$ the law of
$v=\sqrt{U}$ and $\mu:=\nu-p\,\delta_1$,\, $\mu$ is a positive measure
supported in $(0,1)$ with total mass $1-p$, and for every $d\ge 1$ the
exact decomposition \eqref{eq:decomp} holds.
\end{theorem}

\begin{proof}
\emph{Ratios.} For $k\ge 1$, by \eqref{eq:CP},
\[
\log\frac{\E[U^k]}{\E[U^{k-1}]}
=\int_0^\infty\bigl(\e^{-kt}-\e^{-(k-1)t}\bigr)
\,\eta(t)\,\frac{\e^{-t}}{1-\e^{-t}}\,\dd t
=-\int_0^\infty \e^{-kt}\,\eta(t)\,\dd t,
\]
because $\bigl(\e^{-kt}-\e^{-(k-1)t}\bigr)\e^{-t}/(1-\e^{-t})
=-\e^{-kt}$. By Lemma~\ref{lem:binet},
\begin{equation}\label{eq:ratio}
\frac{\E[U^k]}{\E[U^{k-1}]}
=\Bigl(k-\tfrac12\Bigr)\,\e^{\gamma-H_{k-1}} .
\end{equation}
On the other hand, from \eqref{eq:Mk} and the telescoping
$S_1(k)-S_1(k-1)=H_{k-1}$ of \eqref{eq:S1},
\[
\frac{M_k}{M_{k-1}}
=\frac{2k(2k-1)}{4k}\,\e^{\gamma-H_{k-1}}
=\Bigl(k-\tfrac12\Bigr)\,\e^{\gamma-H_{k-1}} .
\]
Since $\E[U^0]=1=M_0$, induction gives $\E[U^k]=M_k$ for all $k$.

\emph{Atom.} Since $U\in(0,1]$, $U^k\downarrow\mathbf 1_{\{U=1\}}$
pointwise, so by dominated convergence
$\Prob(U=1)=\lim_k \E[U^k]=\lim_k M_k$. Stirling's formula gives
\[
\frac{(2k)!}{4^k\,k!}=\sqrt2\,k^k\e^{-k}\,\bigl(1+O(1/k)\bigr),
\]
while $H_{k-1}=\log k+\gamma-\tfrac1{2k}+O(k^{-2})$ and \eqref{eq:S1}
give $\gamma k-S_1(k)=-k\log k+k-\tfrac12+O(1/k)$; multiplying,
$M_k=\sqrt2\,\e^{-1/2}+O(1/k)\to\sqrt{2/\e}$.

\emph{Decomposition.} $v=\sqrt U\in(0,1]$ has an atom of mass exactly
$p$ at $1$ and no other mass on $\{1\}$, so $\mu=\nu-p\delta_1\ge 0$
is supported in $(0,1)$ with mass $1-p$. Finally, for any $d\ge 1$,
expanding the finite sum \eqref{eq:CdPsid} and using
$\E[(\sqrt U)^{2k}]=\E[U^k]=M_k$,
\[
\E\,C_d(\sqrt U\,w)
=\sum_{k=0}^{d}(-1)^k\frac{d^{\underline k}}{d^{\,k}(2k)!}\,
\E[U^k]\;w^{2k}
=\Psi_d(w),
\]
which is \eqref{eq:decomp} after splitting $\nu=p\,\delta_1+\mu$.
\end{proof}

\begin{remark}
Only three consequences of Theorem~\ref{thm:moments} are used below:
$\mu\ge0$, the mass $\mu((0,1))=1-p$, and the support
$\operatorname{supp}\mu\subseteq[0,1]$. The specific shape of
$\lambda$ matters only through the closed form \eqref{eq:Mk} of the
moments. Note also $p^2=2/\e>1/4$, i.e.
\begin{equation}\label{eq:phalf}
p\;>\;\tfrac12 ,
\end{equation}
the single inequality on which the entire sign-transfer argument runs
(it amounts to $8>\e$).
\end{remark}

\section{The critical Laguerre polynomial and its energy}\label{sec:laguerre}

\begin{lemma}[Hermite bridge and zeros]\label{lem:hermite}
Let $\He_n$ denote the probabilists' Hermite polynomials,
$\He_n(x)=(-1)^n\e^{x^2/2}\frac{\dd^n}{\dd x^n}\e^{-x^2/2}$. Then
\begin{equation}\label{eq:bridge}
C_d(w)\;=\;\frac{\He_{2d}\bigl(w/\sqrt{2d}\bigr)}{\He_{2d}(0)},
\qquad
\He_{2d}(0)=(-1)^d\,(2d-1)!! .
\end{equation}
Consequently $C_d$ is an even polynomial of degree $2d$ with
$C_d(0)=1$, and all its zeros are real and simple:
\[
-\rho_d<\dots<-\rho_1<0<\rho_1<\dots<\rho_d .
\]
\end{lemma}

\begin{proof}
The quadratic transformation between Hermite and Laguerre polynomials
\cite[(5.6.1)]{szego}, in confluent-hypergeometric form
\cite[Ch.~18]{dlmf}, reads
$H_{2d}(x)=(-1)^d\frac{(2d)!}{d!}\,\hypF{-d}{\tfrac12}{x^2}$ for the
physicists' polynomials $H_n$; converting by
$\He_n(x)=2^{-n/2}H_n(x/\sqrt2)$ and substituting $x=w/\sqrt{4d}$
yields \eqref{eq:bridge}, the constant being
$\He_{2d}(0)=(-1)^d(2d)!/(2^d d!)=(-1)^d(2d-1)!!$. The zeros of
Hermite polynomials are real and simple \cite[Thm.~3.3.1]{szego}, and
$0$ is not among them for even index since $\He_{2d}(0)\neq 0$;
evenness of $\He_{2d}$ arranges them symmetrically.

\emph{Remark on the formal proof.} In the Lean formalization the
identity \eqref{eq:bridge} is not imported from the literature: both
sides are even polynomials with constant term $1$ whose even
coefficient sequences satisfy the same first-order recurrence
\[
(2k+1)(2k+2)\,a_{k+1}\;=\;\bigl(2k-2d\bigr)\,a_k\,/\,(2d)
\]
forced by their respective differential equations, and the
real-rootedness of $\He_n$, with simplicity and the interlacing used
below, is proved from scratch by induction on $n$ through the
recurrence $\He_{n+1}=x\He_n-\He_n'$ (Section~\ref{sec:lean}).
\end{proof}

\begin{lemma}[Kummer ODE]\label{lem:ode}
$C_d''(w)-\dfrac{w}{2d}\,C_d'(w)+C_d(w)=0$.
\end{lemma}

\begin{proof}
With $u(z)=\hypF{-d}{1/2}{z}$ and $z=w^2/4d$, Kummer's equation
$z\,u''+(\tfrac12-z)\,u'+d\,u=0$ \cite[Ch.~13]{dlmf} transforms under
$C_d(w)=u(w^2/4d)$ by the chain rule:
$C_d'=u'\cdot w/2d$ and
$C_d''=u''\,w^2/4d^2+u'/2d$, so, collecting terms (all derivatives of
$u$ evaluated at $z=w^2/4d$),
\[
C_d''-\frac{w}{2d}C_d'+C_d
=\frac{u''\,w^2}{4d^2}+\frac{u'}{2d}-\frac{u'\,w^2}{4d^2}+u
=\frac{1}{d}\Bigl(z\,u''+\bigl(\tfrac12-z\bigr)u'+d\,u\Bigr)
\;=\;0 . \qedhere
\]
\end{proof}

\begin{lemma}[critical data]\label{lem:critical}
$C_d'$ has exactly $2d-1$ zeros, all real and simple; on $[0,\infty)$
they are
$0=c_0<c_1<\dots<c_{d-1}$, interlacing with the zeros as
\[
0=c_0<\rho_1<c_1<\rho_2<\dots<c_{d-1}<\rho_d ,
\]
and the critical values alternate in sign:
\begin{equation}\label{eq:altsign}
\sgn C_d(c_m)\;=\;(-1)^m,\qquad m=0,\dots,d-1 .
\end{equation}
For $d=1$ the list of positive critical points is empty and
\eqref{eq:altsign} reduces to $C_1(0)=1>0$.
\end{lemma}

\begin{proof}
$C_d'$ has degree $2d-1$; Rolle's theorem places at least one of its
zeros strictly between consecutive zeros of $C_d$, which accounts for
$2d-1$ of them and therefore for all, each simple, one per gap. The
gap $(-\rho_1,\rho_1)$ contains the critical point $0$ (evenness). On
$(\rho_m,\rho_{m+1})$ the polynomial $C_d$ has no zero, hence constant
sign; the sign on $(-\rho_1,\rho_1)$ is $+$ because $C_d(0)=1$, and it
flips across each simple zero $\rho_m$, so the sign on
$(\rho_m,\rho_{m+1})$ is $(-1)^m$. Since
$c_m\in(\rho_m,\rho_{m+1})$ for $1\le m\le d-1$ and $c_0=0$, the
display \eqref{eq:altsign} follows.
\end{proof}

\begin{lemma}[energy monotonicity and segment domination]\label{lem:energy}
Let $E(x)=C_d(x)^2+C_d'(x)^2$. Then $E'(x)=\dfrac{x}{d}\,C_d'(x)^2\ge0$
for $x\ge 0$, and $E(0)=1$. Consequently, for every critical point
$c\ge 0$ of $C_d$ and every $v\in[0,1]$,
\begin{equation}\label{eq:domination}
|C_d(vc)|\;\le\;|C_d(c)| ,
\qquad\text{and}\qquad
|C_d(c)|\;\ge\;1 .
\end{equation}
\end{lemma}

\begin{proof}
Differentiate and substitute the ODE of Lemma~\ref{lem:ode}:
\[
E'=2C_d'\,\bigl(C_d+C_d''\bigr)=2C_d'\cdot\frac{x}{2d}\,C_d'
=\frac{x}{d}\,C_d'^2\;\ge\;0\qquad(x\ge0).
\]
$E(0)=C_d(0)^2+C_d'(0)^2=1+0$. Monotonicity gives, for $0\le v\le1$,
\[
C_d(vc)^2\;\le\;E(vc)\;\le\;E(c)\;=\;C_d(c)^2 ,
\]
the last equality because $C_d'(c)=0$; and
$C_d(c)^2=E(c)\ge E(0)=1$.
\end{proof}

\begin{remark}
Lemma~\ref{lem:energy} is the real-axis case of an energy monotonicity
that in fact holds along every horizontal line in $\C$
($\partial_x(|C_d|^2+|C_d'|^2)=(x/d)|C_d'|^2$, the imaginary part of
$w$ contributing no real part to the cross terms). Only the real-axis
case is needed here, and only it was formalized. We also note that
\eqref{eq:domination} with constant $1$ is what makes the argument
work with no room to spare; see Remark~\ref{rem:sharp}.
\end{remark}

\section{Proof of the main theorem}\label{sec:proof}

\begin{proof}[Proof of Theorem~\ref{thm:main}]
Fix $d\ge1$ and abbreviate $p=\sqrt{2/\e}$. Recall from
\eqref{eq:decomp} that
$\Psi_d(x)=p\,C_d(x)+\int_{(0,1)}C_d(vx)\,\dd\mu(v)$ with $\mu\ge0$ of
total mass $1-p$.

\emph{Step 1: $\Psi_d$ copies the critical signs of $C_d$.} Let
$c\in\{c_0,\dots,c_{d-1}\}$ be a nonnegative critical point of $C_d$.
By the segment domination \eqref{eq:domination},
\[
\Bigl|\Psi_d(c)-p\,C_d(c)\Bigr|
\;\le\;\int_{(0,1)}\bigl|C_d(vc)\bigr|\,\dd\mu(v)
\;\le\;(1-p)\,\bigl|C_d(c)\bigr|
\;<\;p\,\bigl|C_d(c)\bigr| ,
\]
the last (strict) inequality by \eqref{eq:phalf} and
$|C_d(c)|\ge1>0$. Hence $\Psi_d(c)\ne0$ and
\begin{equation}\label{eq:signcopy}
\sgn\Psi_d(c_m)\;=\;\sgn C_d(c_m)\;=\;(-1)^m,
\qquad m=0,\dots,d-1 .
\end{equation}

\emph{Step 2: tail sign.} The leading coefficient of $\Psi_d$ is
$(-1)^d\,d!\,M_d/\bigl(d^{\,d}(2d)!\bigr)$, of sign $(-1)^d$ since
$M_d>0$; the degree is exactly $2d$. Hence
$\sgn\Psi_d(x)=(-1)^d$ for all sufficiently large $x>0$.

\emph{Step 3: count.} By \eqref{eq:signcopy} and Step~2 the sequence
of values
\[
\Psi_d(c_0),\,\Psi_d(c_1),\,\dots,\,\Psi_d(c_{d-1}),\,
\Psi_d(x_\infty)\qquad(x_\infty\ \text{large})
\]
has signs $+,-,+,\dots,(-1)^{d-1},(-1)^d$: it alternates through $d$
strict sign changes. By the intermediate value theorem $\Psi_d$ has at
least $d$ distinct zeros in $(c_0,\infty)=(0,\infty)$. Since $\Psi_d$
is even, it has at least $d$ further distinct zeros in $(-\infty,0)$,
and $\Psi_d(0)=M_0=1\neq0$. That is at least $2d$ distinct complex
zeros of a polynomial of degree $2d$: there are exactly $2d$, all
real, all simple.
\end{proof}

\begin{remark}[$d=1$]
The degenerate case is instructive: there are no positive critical
points, the sign sequence is $\Psi_1(0)=1>0$ followed by the tail sign
$-$, and indeed
$\Psi_1(w)=1-\tfrac{\e^{\gamma}}{4}w^2$ has the two real simple zeros
$\pm2\,\e^{-\gamma/2}$. The proof above covers this case verbatim ---
no step assumed $d\ge2$.
\end{remark}

\section{The formal proof}\label{sec:lean}

Theorem~\ref{thm:main} (in the realness form stated in
Section~\ref{sec:intro}) has been formalized in Lean~4 (v4.30.0)
against mathlib (pinned at commit
\texttt{c5ea0035}\,/\,tag v4.30.0). The development comprises
thirteen files, about $5{,}300$ lines:

\begin{center}
\footnotesize
\begin{tabular}{@{}l p{0.64\textwidth}@{}}
\texttt{Defs} & the polynomials $\Psi_d$, $C_d$, the sequences $M_k$, $S_1$\\
\texttt{Structure} & coefficient closed forms, degree, the Kummer ODE in polynomial form\\
\texttt{Frullani} & Frullani's integral for exponentials\\
\texttt{Binet} & the integer Binet identity (Lemma~\ref{lem:binet})\\
\texttt{CPMeasure} & the L\'evy measure $\lambda$ as a \texttt{withDensity} measure; finiteness\\
\texttt{Moments}, \texttt{MomentsLimit} & moment ratios; the Stirling limit $M_k\to p$\\
\texttt{MuBridge} & the finite-quadrature interface: residual moments $M_k-p$ from a measure\\
\texttt{Capstone} & discharge of the quadrature interface by the compound-Poisson pushforward\\
\texttt{Energy} & Lemma~\ref{lem:energy} (real energy monotonicity, segment domination)\\
\texttt{Hermite} & real-rootedness, simplicity and interlacing of $\He_n$, by induction\\
\texttt{CriticalData} & Lemma~\ref{lem:critical} for $C_d$ via the bridge \eqref{eq:bridge}; statement and proof of \texttt{theorem\_M} and \texttt{theorem\_M\_aeval}\\
\texttt{SignCount} & the sign-alternation count, Step~3 above\\
\end{tabular}
\end{center}

Two design points deserve mention.

\emph{Measures, not samples.} The formal development verifies every
measure-level assertion of Section~\ref{sec:moments}: positivity and
finiteness of $\lambda$ (a \texttt{withDensity} measure), the integer
Binet identity as a genuine Lebesgue-integral statement, and the
existence of the residual measure $\mu$ of \eqref{eq:decomp} ---
obtained as a pushforward --- with $\mu\ge0$, mass $1-p$ and even
moments $M_k-p$. What it does not do is represent $U$ as the
exponential functional $\e^{-\sum_j T_j}$ of a Poisson point process.
This is a deliberate boundary, not a gap: the proof of
Theorem~\ref{thm:main} consumes the law of $U$ only through the
listed measure data, so the point-process representation --- the right
way to \emph{understand} the law, and the way it was found
(Section~\ref{sec:moments}) --- contributes no additional logical
content to the theorem. Formalizing it would enlarge the verified
statement about the measure, not the verified theorem about
$\Psi_d$.

\emph{Hermite theory from scratch.} Rather than formalize the
classical literature on Laguerre zeros, the development proves
real-rootedness with simplicity and the full interlacing data for the
probabilists' Hermite polynomials by induction on $n$ via
$\He_{n+1}=x\,\He_n-\He_n'$: between consecutive simple roots of
$\He_n$ the derivative term forces a sign change of $\He_{n+1}$, and
two tail roots are produced by the parity of the leading terms; a
counting argument then pins exactly $n+1$ simple roots. The identity
\eqref{eq:bridge} transports this package to $C_d$ by a first-order
coefficient recurrence, avoiding any appeal to unformalized special
function theory.

The final check reports, for the headline theorem and its
evaluation-form companion alike,
\begin{center}
\texttt{'TheoremM.theorem\_M' depends on axioms:
[propext, Classical.choice, Quot.sound]}\\
\texttt{'TheoremM.theorem\_M\_aeval' depends on axioms:
[propext, Classical.choice, Quot.sound]}
\end{center}
with zero \texttt{sorry}s in the tree (8489 compiled jobs). These
three axioms are the standard foundation of all of mathlib;
in particular no numerical computation, floating point or interval
arithmetic enters the formal proof anywhere --- all constants,
including $p=\sqrt{2/\e}$ and the inequality \eqref{eq:phalf}, are
handled exactly.

The formalization was produced jointly by two AI research agents
(Claude ``Fable'' and GPT) working under the author's direction in a
mutual-audit protocol (disjoint file ownership, cross line-audits of
each other's proofs), with the Lean kernel as the final referee. In
addition, the development has been re-verified independently of
Lean's own kernel: the proof terms were exported with the upstream
\texttt{lean4export} tool (v4.30.0) and re-checked by Nanoda
(\texttt{nanoda\_lib}, commit \texttt{f58f2f6d}), an external
implementation of the Lean~4 kernel, which reports
\emph{``Checked 48241 declarations with no errors''} for an export
containing \texttt{theorem\_M} and \texttt{theorem\_M\_aeval}, with
the permitted axioms restricted to exactly the three listed above.

\section{Concluding remarks}\label{sec:remarks}

\begin{remark}[sharpness]\label{rem:sharp}
The budget of Step~1 is sharp at the endpoint $v\to1^-$: numerically,
for every $d$ tested ($1\le d\le 300$), the worst ratio
$|\Psi_d(c_m)-pC_d(c_m)|\,/\,p|C_d(c_m)|$ over the critical points
equals $(1-p)/p=0.16582\ldots$ to all computed digits --- the measure
$\mu$ concentrates its action at $v$ near $1$, where the domination
\eqref{eq:domination} costs nothing. The proof therefore has a uniform
factor-$6$ margin, but no asymptotic slack is hiding anywhere: the
mechanism is the strict inequality $p>\tfrac12$, used once.
\end{remark}

\begin{remark}[simplicity vs.\ the formal statement]
The paper proof gives realness \emph{and simplicity} (Step~3 produces
$2d$ \emph{distinct} zeros). The formalized statement asserts realness
of every complex root; distinctness is present in the formal
development's counting layer but was not exported into the headline
statement.
\end{remark}

\begin{remark}[a question]
The moment sequence \eqref{eq:Mk} acts coefficientwise on $C_d$ and
preserves hyperbolicity for every $d$. Is $(M_k)_{k\ge0}$ a
P\'olya--Schur multiplier sequence (i.e.\ does it preserve
hyperbolicity of \emph{all} hyperbolic even polynomials)? The
compound-Poisson representation proves much less --- it is tied to the
energy structure of $C_d$ --- and we do not know the answer; a
negative answer would make Theorem~\ref{thm:main} more interesting,
not less.
\end{remark}

\subsection*{Acknowledgments}
% TODO (Michael): authorship/credit framing is yours to settle.
The mathematical route and its formal verification were developed
jointly by two AI research agents, Claude (``Fable'', Anthropic) and
GPT (OpenAI), working as a coordinated pair under the author's
direction, with mutual line-audits and the Lean kernel as referee.
Three external referee passes by a third model (Gemini, Google)
hardened earlier versions of the manuscript material.

\begin{thebibliography}{9}

\bibitem{szego}
G.~Szeg\H{o},
\emph{Orthogonal Polynomials},
4th ed., Amer. Math. Soc. Colloq. Publ., vol.~23,
American Mathematical Society, Providence, RI, 1975.

\bibitem{polya-schur}
G.~P\'olya and J.~Schur,
\emph{\"Uber zwei Arten von Faktorenfolgen in der Theorie der
algebraischen Gleichungen},
J. Reine Angew. Math. \textbf{144} (1914), 89--113.

\bibitem{dlmf}
NIST Digital Library of Mathematical Functions,
\url{https://dlmf.nist.gov/},
Release 1.2.x; F.~W.~J. Olver et al., eds.

\bibitem{ww}
E.~T. Whittaker and G.~N. Watson,
\emph{A Course of Modern Analysis},
4th ed., Cambridge University Press, 1927.

\bibitem{kingman}
J.~F.~C. Kingman,
\emph{Poisson Processes},
Oxford Studies in Probability, vol.~3,
Oxford University Press, 1993.

\bibitem{lean4}
L.~de~Moura and S.~Ullrich,
\emph{The Lean 4 theorem prover and programming language},
in: Automated Deduction -- CADE 28,
Lecture Notes in Comput. Sci., vol.~12699, Springer, 2021, 625--635.

\bibitem{mathlib}
The mathlib Community,
\emph{The Lean mathematical library},
in: Proc. 9th ACM SIGPLAN Int. Conf. on Certified Programs and Proofs
(CPP 2020), ACM, 2020, 367--381.

\end{thebibliography}

\end{document}
